% ==================================================
% Falling Factorials
% Discrete Calculus --- Polynomial Basis
% ==================================================

\subsection{Polynomial Basis: Falling Factorials}
\label{sec:falling-factorials}

\begin{exposition}
In ordinary calculus, polynomial functions are naturally expressed using
the monomial powers \(1, x, x^2, x^3, \dots\) because the derivative
acts simply on them: \(\tfrac{d}{dx} x^n = n x^{n-1}\).

In discrete calculus, however, the difference operator does not behave
as cleanly on ordinary powers. For example,
\[
\Delta(n^2) = (n+1)^2 - n^2 = 2n + 1.
\]
Although the degree decreases by one, the result contains an extra
constant term. For this reason it is more natural to work with a
different family of polynomials --- the \emph{falling factorials} ---
which interact with \(\Delta\) exactly as monomials interact with
\(\tfrac{d}{dx}\).
\end{exposition}

\begin{definitionbox}{Definition (Falling Factorial)}
\begin{definition}[Falling Factorial]
\label{def:falling-factorial}
For a nonnegative integer \(k\), the \emph{falling factorial} (or
\emph{falling power}) is defined by
\[
x^{\underline{k}} = x(x-1)(x-2)\cdots(x-k+1)
\qquad (k \text{ factors}).
\]
By convention, \(x^{\underline{0}} = 1\).
\end{definition}
\end{definitionbox}

\begin{remark*}[Standard quantified statement]
\(\forall k \in \mathbb{N},\; \forall x \in \mathbb{R}:\;
x^{\underline{k}} := \displaystyle\prod_{j=0}^{k-1}(x-j)\),
with \(x^{\underline{0}} := 1\).
\end{remark*}

\begin{remark*}[Predicate reading]
The expression \(x^{\underline{k}}\) denotes the product of the \(k\)
successive descending factors beginning at \(x\), with the zero-factor case
defined to be \(1\).
\end{remark*}

\begin{remark*}[Interpretation]
When \(x = n \in \mathbb{N}\) and \(k \le n\), the falling factorial
equals \(\tfrac{n!}{(n-k)!}\). In particular, \(n^{\underline{n}} = n!\).
Thus falling factorials generalise factorials to arbitrary inputs.
\end{remark*}

\DefinitionalRoot

\begin{example*}[Examples]
The first few falling factorials are:
\[
\begin{aligned}
x^{\underline{0}} &= 1, \\
x^{\underline{1}} &= x, \\
x^{\underline{2}} &= x(x-1), \\
x^{\underline{3}} &= x(x-1)(x-2).
\end{aligned}
\]
\end{example*}

\subsection*{The Difference Rule for Falling Factorials}

\begin{theorembox}{Theorem (Difference Rule for Falling Factorials)}
\begin{theorem}[Difference Rule for Falling Factorials]
\label{thm:falling-factorial-rule}
For \(k \ge 1\),
\[
\Delta\, x^{\underline{k}} = k\,x^{\underline{k-1}}.
\]
\hyperref[prf:falling-factorial-rule]{\textit{Go to proof.}}
\end{theorem}
\end{theorembox}

\begin{remark*}[Standard quantified statement]
\(\forall k \in \mathbb{N},\; k \ge 1,\; \forall x \in \mathbb{R}:\;
(x+1)^{\underline{k}} - x^{\underline{k}} = k\,x^{\underline{k-1}}\).
\end{remark*}

\begin{remark*}[Predicate reading]
\[
\Delta x^{\underline{k}}=k x^{\underline{k-1}}.
\]
\end{remark*}

\begin{remark*}[Interpretation]
This formula exactly mirrors the continuous power rule
\(\tfrac{d}{dx}x^k = kx^{k-1}\). Because of this clean behaviour under
\(\Delta\), falling factorials form the natural polynomial basis for
discrete calculus.
The theorem also yields the discrete antiderivative rule:
\[
\Delta\!\left(\frac{x^{\underline{k+1}}}{k+1}\right) = x^{\underline{k}}.
\]
Thus \(\tfrac{x^{\underline{k+1}}}{k+1}\) is a discrete antiderivative
of \(x^{\underline{k}}\), mirroring the continuous rule
\(\int x^k\,dx = \tfrac{x^{k+1}}{k+1}\).
\end{remark*}

\begin{dependencies}
\begin{itemize}
  \item \hyperref[def:falling-factorial]{Falling Factorial}
  \item \hyperref[def:forward-difference]{Forward Difference Operator}
\end{itemize}
\end{dependencies}

\subsection*{Summary Table}

\begin{exposition}
\begin{center}
\renewcommand{\arraystretch}{1.4}
\begin{tabular}{c|c}
\(f(x)\) & \(\Delta f(x)\) \\
\hline
\(x^{\underline{1}} = x\) & \(1\) \\
\(x^{\underline{2}} = x(x-1)\) & \(2x\) \\
\(x^{\underline{3}} = x(x-1)(x-2)\) & \(3x(x-1)\) \\
\(x^{\underline{4}} = x(x-1)(x-2)(x-3)\) & \(4x(x-1)(x-2)\)
\end{tabular}
\end{center}
\end{exposition}

\subsection*{Polynomial Representation in the Falling Factorial Basis}

\begin{exposition}
Every polynomial can be written as a linear combination of falling
factorials:
\[
f(x)
=
c_0
+ c_1\, x
+ c_2\, x^{\underline{2}}
+ c_3\, x^{\underline{3}}
+ \cdots
+ c_n\, x^{\underline{n}}.
\]
\end{exposition}

\begin{remark*}[Interpretation]
This representation plays the same role in discrete calculus that the
Taylor expansion plays in continuous calculus. The coefficients \(c_k\)
are determined by the successive finite differences of \(f\) evaluated
at \(0\): \(c_k = \Delta^k f(0) / k!\). This leads directly to Newton's
forward difference expansion, developed in the Expansions section.
\end{remark*}
